
本稿は、意識を内部で生成される対象としてではなく、異なる時間スケールをもつ処理のあいだに一時的に開かれる状態として定式化する。中心となる主張は、この記述層において意識を特定の脳部位、表象、あるいは出力として扱うのではなく、高速な予測過程と低速な沈殿的記憶過程のあいだで成立するログ参照可能な関係として扱うという点にある。

この関係を記述するために、内部遅延またはバッファリングを $\tau$、有効なトルク様量を $\mathcal{T}$、情報ポテンシャルを $\Phi$ とする。速度差は $\Delta v = v_f - v_s$、結合粘性は $\mu$ と表す。Libet 型実験、分離脳実験、クオリア、測定、非局在性はいずれも、この速度スケール間の変換と閉包不全の観点から再解釈される。

さらに本稿は、ログ参照を意味勾配に沿った非閉包的な情報流として再記述する。その際、IFGT の情報密度 $I$、情報流 $J$、情報ポテンシャル $\Phi$ は、一般情報場そのものと意識を直同一視するためではなく、認知的観測窓に射影された有効変数として用いられる。
